\section{Just log-concavity}
\label{sec:log_factor_reduction_to_tail}
The next result shows that if $q$ is Laplace in dimension one, then we can still get the desired $\tilde{O}(\eta^2)$ bound by paying a logarithmic factor.
This also shows that $f_\alpha$-divergence is crucial and one would lose a lot (literally infinite)
if simply using the $\chi^2$-divergence as a surrogate for $f_\alpha$-divergence.
We prove this through a sequence of lemmas that reduce the problem of bounding $D_{f_\alpha}(\PP_{X, T_0} || \PP_X \otimes \PP_{T_0})$ to bounding the tail probability of $R_0$.
Recall
\begin{align}
    R_0=\frac{P_{X,T_0}(X, T_0)}{P_X(X)P_{T}(T_0)}\,, \qquad
    \text{and}\qquad
    R=\frac{P_{X,T_0}(X, T)}{P_X(X)P_T(T)} \qquad \text{for } T \stackrel{i.i.d}{\sim} p_{T_0}\,,
\end{align}

\begin{lemma}[$f_\alpha$ is qaudratic-then-linear]
    \label{lem:log_factor_fa_quadratic_then_linear}
    For every $\alpha\in(0,1)$ and $P\ll Q$ we have
    \begin{align}
        D_{f_\alpha}(P\|Q)
         & \le
        e^{\alpha}
        \alpha^2
        \left(
        1+\E[R_0\mathbf{1}_{\{R_0\le \frac{1}{\alpha}\}}]
        \right)
        +
        \alpha \PP\left(R_0>\frac{1}{\alpha}\right)\,.
        \label{eq:log_factor_reduction_to_tail}
    \end{align}
\end{lemma}
The proof is deferred to \cref{proof:log_factor_fa_quadratic_then_linear}.

The next result shows that how a tail bound on $R_0$ implies a bound on $D_{f_\alpha}(\PP_{X, T_0} || \PP_X \otimes \PP_{T_0})$.
\begin{lemma}[Tail to divergence control]
    \label{lem:log_factor_tail_to_divergence}
    Suppose that for every $y\ge1$,
    \marginpar{\textit{Note that if we had $\E[R_0] = C <\infty$, then this bound would be trivial.
            But $\E[R_0]$ is the chi-squared divergence, which is infinite.}}
    \begin{align}
        \PP(R_0>y)
        \le
        \frac{C}{y}\,,
        \label{eq:log_factor_R0_tail}
    \end{align}
    for some constant $C>0$.
    Then for every $\alpha\in(0,1)$,
    \begin{align}
        D_{f_\alpha}(P\|Q)
         & \le
        C \alpha^2\log\left(\frac{e}{\alpha}\right).
        \label{eq:log_factor_fa_upper}
    \end{align}
\end{lemma}
The proof is deferred to \cref{proof:log_factor_tail_to_divergence}.

\marginpar{Basically we are almost sure that $R_0$ has a Pareto tail.}
Now, it is clear that if we can show the tail bound in \eqref{eq:log_factor_R0_tail} for $R_0$, then we get the desired $\tilde{O}(\eta^2)$ bound on the variance term.
We first do it for the Laplace distribution, and then we show that the same tail bound holds for a general log-concave distribution, in 1 dimension.

\noindent
First, define
\begin{align}
    X'\triangleq\lambda\sum_{k=1}^\infty (1-\lambda)^{k-1} T_k,
    \qquad
    \text{so that }\qquad
    X=\lambda T_0+(1-\lambda) X'.
\end{align}
Note that $X' \stackrel{d.}{\equiv} X$ and $X' \perp T_0$.
Let $T \stackrel{i.i.d.}{\sim} T_0$; the following identity for $R_0$ plays a crucial role in the analysis:
\begin{align}
    R_0
     & =
    \frac{p_{X|T_0}(X)}{p_X(X)}
    =
    \frac{p_{X|T_0}(\lambda T_0 + (1-\lambda) X')}{p_{X}(X)}
    =
    \frac{1}{(1-\lambda)^d}\cdot \frac{p_X(X')}{p_X(X)}
    \label{eq:R0_identity_1}
    \\
     & =
    \frac{1}{(1-\lambda)^d}\cdot \frac{p_X(X')}{p_X(\lambda T_0+(1-\lambda)X')}
    =
    \frac{1}{(1-\lambda)^d}\cdot \frac{p_X(X)}{p_X(\lambda T+(1-\lambda)X)}\,.
    \label{eq:R0_identity_2}
\end{align}

\subsubsection{Bounds in 1d}
We now show that how the tail bounds on $R_0$ can be established in the 1-dimensional case when $q$ is Laplace, and then we show that the same tail bound holds for a general log-concave distribution in 1 dimension.

\begin{theorem}[Laplace bound]
    \label{theorem:laplace_log_factor}
    Assume $d=1$, and let $q(t)=\frac12 e^{-|t|}$.
    Then there exists a numerical constant $C>0$ such that for every
    $\lambda\in(0,1)$ and every $\alpha\in(0,1]$,
    \begin{align}
        D_{f_\alpha}(\PP_{X,T_0}\|\PP_X\otimes\PP_{T_0})
        \le
        \frac{4}{1-\lambda}
        \alpha^2 \log\left(\frac{e}{\alpha}\right)\,.
        \label{eq:laplace_log_factor_fa_upper}
    \end{align}
    % In particular,
    % \begin{align}
    %     \sup_{f:\K\to[0,1]}
    %     \Lambda_{\mathrm{var},q}(E_\lambda,p_X;f)
    %     \le
    %     \frac{C}{1-\lambda}
    %     \frac{\eta^2}{\lambda^2}
    %     \log\left(\frac{e\lambda}{\eta}\right)
    %     +
    %     \frac{\eta^2}{2}.
    %     \label{eq:laplace_log_factor_variance_upper}
    % \end{align}
    % Moreover, the logarithm is not an artefact: there exist constants
    % $c_\lambda>0$ and $a_\lambda\in(0,1]$, depending only on $\lambda$,
    % such that for every $a\in(0,a_\lambda]$,
    % \begin{align}
    %     D_{f_a}(\PP_{X,T_0}\|\PP_X\otimes\PP_{T_0})
    %     \ge
    %     c_\lambda a^2 \log\left(\frac1a\right).
    %     \label{eq:laplace_log_factor_fa_lower}
    % \end{align}
\end{theorem}
% The proof is deferred to \cref{proof:laplace_log_factor}.
% \subsubsection{Proof of \cref{theorem:laplace_log_factor}}
% \label{proof:laplace_log_factor}
\begin{proof}
    For a 1-dimensional application of \cref{eq:R0_identity_2}, we have
    \begin{align}
        R_0
         & =
        \frac{1}{1-\lambda}\cdot \frac{p_X(X)}{p_X(\lambda T+(1-\lambda)X)}\,.
        \label{eq:laplace_log_factor_R0_formula}
    \end{align}
    Next, we will show that $p_X$ is log-Lipschitz with Lipschitz constant $1/\lambda$.
    Let
    \begin{align}
        q_\lambda(u)=\frac{1}{2\lambda}e^{-|u|/\lambda}
    \end{align}
    be the density of $\lambda T_0$.
    Let $p_z$ be the density of $Z = (1-\lambda)X'$,
    then since $p=q_\lambda*p_z$, we have the lower bound
    \marginpar{
        \textit{There is an upper bound too:}
        \begin{align}
            \frac{p(x+h)}{p(x)}
            \le
            e^{|h|/\lambda}\,.
        \end{align}
    }
    \begin{align}
        p(x+h)
         & =
        \int_\R q_\lambda(x+h-z)p_z(z)\,dz
        \\
         & \ge
        e^{-|h|/\lambda}
        \int_\R q_\lambda(x-z)p_z(z)\,dz
        =
        e^{-|h|/\lambda}p(x)\,,
        \label{eq:laplace_log_factor_log_lipschitz}
    \end{align}
    which implies that $p$ is $\frac{1}{\lambda}$-log-Lipschitz.
    On the other hand, using the same self-identical decomposition,
    \begin{align}
        p(\lambda t+(1-\lambda)z)
         & =
        \frac{1}{1-\lambda}
        \int_\R
        q(s)
        p\left(
        z+\frac{\lambda(t-s)}{1-\lambda}
        \right)
        ds
        \\
         & \ge
        \frac{p(z)}{1-\lambda}
        \int_\R
        q(s)
        \exp\left(
        -\frac{|t-s|}{1-\lambda}
        \right)
        ds,
        \label{eq:laplace_log_factor_denominator_lower}
    \end{align}
    where we used \eqref{eq:laplace_log_factor_log_lipschitz}.
    Thus, we have
    \begin{align}
        R_0
        =
        \frac{1}{1-\lambda}\cdot \frac{p_X(X)}{p_X(\lambda T+(1-\lambda)X)}
         & \le
        \left(
        \int_\R
        q(s)
        \exp\left(
        -\frac{|T-s|}{1-\lambda}
        \right)
        \d s
        \right)^{-1}
        \label{eq:laplace_log_factor_R0_integral_lower}
        \\
         & =
        \E\left[
            \exp\left(
            -\frac{|T - T'|}{1-\lambda}
            \right)
            \middle | T
            \right]^{-1}\,.
        \label{eq:laplace_log_factor_R0_integral_upper}
    \end{align}
    Since $q(s)=\frac12e^{-|s|}$, for every $t\in\R$,
    \begin{align}
        \E\left[
            \exp\left(
            -\frac{|T - T'|}{1-\lambda}
            \right)
            \middle | T
            \right]
         & =
        \int_\R
        q(s)
        \exp\left(
        -\frac{|T-s|}{1-\lambda}
        \right)
        \d s
        \\
         & =
        2(1-\lambda)
        \int_\R
        q(s)
        \ell_{1-\lambda}(T-s)
        \d s
        \\
         & =
        2(1-\lambda) q * \ell_{1-\lambda}(T)
        \\
         & =
        2(1-\lambda) p_{U}(T)
        \label{eq:laplace_log_factor_laplace_convolution_lower}
    \end{align}
    where $\ell_{1-\lambda}(v)$ is the density of $(1-\lambda)L$ where $L\sim\mathrm{Lap}(0,1)$, and $p_U$ is the density of $U = T + (1-\lambda)L$.
    So we have
    \begin{align}
        R_0
         & \le
        \frac{1}{2(1-\lambda) p_U(T)}
        \quad
        \Rightarrow
        \quad
        \PP(R_0>y)
        \le
        \PP\left(p_U(T)<\frac{1}{2(1-\lambda)y}\right).
    \end{align}

    \marginpar{All the steps so far hold for a general \textbf{globally log-Lipschitz} density $q$ through rescaling.}
    \paragraph{Laplace specific computation.}
    When $t\ge0$, integration gives
    \begin{align}
        \int_\R
        q(s)
        \exp\left(
        -\frac{|t-s|}{1-\lambda}
        \right)
        \d s
         & =
        \frac{1}{2}
        \int_\R
        \exp\left(
        -|s|
        -\frac{|t-s|}{1-\lambda}
        \right)
        \d s    \\
         & \geq
        \frac{1}{2}
        \int_t^\infty
        \exp\left(
        -s
        -\frac{s-t}{1-\lambda}
        \right)
        \d s    \\
         & =
        \frac{1}{2}
        \int_t^{\infty}
        \exp\left(
        -\frac{(2- \lambda) s - t}{1-\lambda}
        \right)
        \d s
        \\
         & =
        \frac{1}{2}
        \int_{0}^\infty
        \exp\left(
        -\frac{(2 - \lambda) (t+s) - t}{1-\lambda}
        \right)
        \d s
        \\
         & =
        \frac{1}{2}
        \int_{0}^\infty
        \exp\left(
        -t - \frac{(2-\lambda)s}{1-\lambda}
        \right)
        \d s
        \\
         & =
        \frac{1}{2}
        \exp(-t)
        \int_{0}^\infty
        \exp\left(
        - \frac{2-\lambda}{1-\lambda} s
        \right)
        \d s
        \\
         & =
        \frac{1-\lambda}{4-2\lambda}
        \exp(-t)\,,
    \end{align}
    and the case $t<0$ is symmetric.
    Combining the last display with \cref{eq:laplace_log_factor_R0_integral_lower} gives
    \begin{align}
        R_0
        \le
        \left(
        \int_\R q(s) \exp\left(- \frac{|T-s|}{1-\lambda}\right) \d s
        \right)^{-1}
        \le
        \frac{4-2\lambda}{1-\lambda}e^{|T|}
        \le
        \frac{4}{1-\lambda}e^{|T|}
        \,.
        \label{eq:laplace_log_factor_R0_exp_upper}
    \end{align}
    Finally, $|T|\sim\mathrm{Exp}(1)$, so
    \begin{align}
        \PP\left(R_0>y\right)
         & \le
        \PP\left(e^{|T|}>\frac{1-\lambda}{4}y\right)\le \frac{4}{(1-\lambda)y}.
    \end{align}
    The result follows from \cref{lem:log_factor_tail_to_divergence}.
\end{proof}

\begin{theorem}[Log-concave kernels pay at most one logarithm]
    \label{theorem:log_concave_log_factor}
    Assume $d=1$, and let $q$ be a log-concave probability density on an interval of $\R$.
    Then there exists a numerical constant $C>0$ such that for every
    $\lambda\in(0,1)$ and every $a\in(0,1]$,
    \begin{align}
        D_{f_a}(\PP_{X,T_0}\|\PP_X\otimes\PP_{T_0})
        \le
        \frac{C}{1-\lambda}
        a^2 \log\left(\frac{e}{a}\right).
        \label{eq:log_concave_log_factor_fa_upper}
    \end{align}
    In particular, for every $0<\eta\le \lambda$,
    \begin{align}
        \sup_{f:\K\to[0,1]}
        \Lambda_{\mathrm{var},q}(E_\lambda,p_X;f)
        \le
        \frac{C}{1-\lambda}
        \frac{\eta^2}{\lambda^2}
        \log\left(\frac{e\lambda}{\eta}\right)
        +
        \frac{\eta^2}{2}.
        \label{eq:log_concave_log_factor_variance_upper}
    \end{align}
\end{theorem}
% The proof is deferred to \cref{proof:log_concave_log_factor}.
% \subsubsection{Proof of \cref{theorem:log_concave_log_factor}}
% \label{proof:log_concave_log_factor}
\begin{proof}
    We prove that
    \begin{align}
        \PP(R_0>y)
        \le
        \frac{C}{(1-\lambda)y},
        \qquad y>0.
        \label{eq:log_concave_log_factor_R0_tail}
    \end{align}
    Let $p=p_X$ be the density of $X$.
    Since $q$ is log-concave and log-concavity is preserved by affine transformations,
    convolution, and weak limits, $p$ is also log-concave.

    Define
    \begin{align}
        X'=\lambda\sum_{k=1}^\infty (1-\lambda)^{k-1} T_k,
        \qquad
        \text{so that }\qquad
        X=\lambda T_0+(1-\lambda) X'.
    \end{align}
    Note that $X' \stackrel{d.}{\equiv} X$ and $X' \perp T_0$,
    so we have
    \begin{align}
        R_0
        =
        \frac{p_{X|T_0}(X)}{p_X(X)}
         & =
        \frac{p_{X|T_0}(\lambda T_0 + (1-\lambda) X')}{p_{X}(X)}
        =
        \frac{1}{1-\lambda}\cdot \frac{p_X(X')}{p_X(X)}\,,\label{eq:log_concave_log_factor_R0_formula}
    \end{align}
    Let $\|p\|_\infty = \sup_{x\in\R} p(x)$.
    From \eqref{eq:log_concave_log_factor_R0_formula}, we get the pointwise bound
    \marginpar{\textit{This inequality is essentially very crude, and potentially should be improved to get the right rate in higher dimensions.}}
    \begin{align}
        R_0
        \le
        \frac{1}{1-\lambda}\frac{\|p\|_\infty}{p(X)}.
        \label{eq:log_concave_log_factor_R0_density_upper}
    \end{align}
    Thus it remains to control how often a one-dimensional log-concave density is much smaller than its maximum.
    We use the following consequence of log-concavity.
    \begin{lemma}[1D Log-concave density concentration]
        Let $Y$ have log-concave density $h$, and let $\|h\|_\infty=h(m)$ for a mode $m$.
        Then, for every $\epsilon > 0$,
        \begin{align}
            \PP\left(h(Y)\le \epsilon \|h\|_\infty\right)
            \le
            2\epsilon.
            \label{eq:log_concave_density_small_ball}
        \end{align}
    \end{lemma}
    Applying \eqref{eq:log_concave_density_small_ball} with $h=p$ and
    \begin{align}
        \epsilon=\frac{1}{(1-\lambda)t},
    \end{align}
    we obtain
    \begin{align}
        \PP(R_0>t)
         & \le
        \PP\left(
        p(X)<\frac{\|p\|_\infty}{(1-\lambda)t}
        \right) \le
        \min\left\{
        1,
        \frac{2}{(1-\lambda)t}
        \right\}.
    \end{align}
    This proves \eqref{eq:log_concave_log_factor_R0_tail}.
\end{proof}

\subsection{Log-concave measures in higher dimensions}
\begin{lemma}[Truncated exponential moment of a log-concave density]
    \label{lem:log_concave_truncated_exp_potential}
    Let $h$ be a log-concave probability density on $\Rd$. Write
    \begin{align}
        h(x)=\|h\|_\infty e^{-V(x)},
    \end{align}
    where $V:\Rd\to[0,\infty]$ is convex and $\inf_x V(x)=0$.
    Let $Y\sim h$. Then there exists a numerical constant $C>0$ such that
    for every $L\ge0$,
    \begin{align}
        \E\left[
            \min\left\{\frac{\|h\|_\infty}{h(Y)},e^L\right\}
            \right]
        =
        \E\left[
            \min\left\{e^{V(Y)},e^L\right\}
            \right]
        \le
        C
        \sum_{j=0}^d \frac{L^j}{j!}.
        \label{eq:log_concave_truncated_exp_potential}
    \end{align}
\end{lemma}

\begin{proof}
    Let
    \begin{align}
        S_d(L)=\sum_{j=0}^d\frac{L^j}{j!}.
    \end{align}
    The result is trivial for $L\le d$, since
    \begin{align}
        \E\left[
            \min\left\{e^{V(Y)},e^L\right\}
            \right]
        \le e^L
        \le 2 S_d(L),
    \end{align}
    where the last inequality is the standard Poisson-tail estimate
    $e^L\le 2\sum_{j=0}^d L^j/j!$ for $L\le d$.

    We now consider $L>d$.
    After translating, assume that the minimum of $V$ is attained at $0$.
    For $s\ge0$, define
    \begin{align}
        K_s=\{x\in\Rd:V(x)\le s\},
        \qquad
        F(s)=\mathrm{vol}(K_s).
    \end{align}
    By convexity of $V$, for every $0<s\le t$,
    \begin{align}
        \frac{s}{t}K_t\subseteq K_s.
    \end{align}
    Hence
    \begin{align}
        \frac{F(s)}{s^d}
        \ge
        \frac{F(t)}{t^d},
        \qquad 0<s\le t.
        \label{eq:sublevel_volume_ratio_monotone}
    \end{align}
    First, by integration by parts, we have
    \begin{align}
        1
        \geq
        \int_{K_L} h(x)\d x
         & =
        \|h\|_\infty
        \int_{V(x)\le L} e^{-V(x)}\d x \\
         & =
        \|h\|_\infty
        \int_0^L e^{-s}\d F(s)         \\
         & =
        \|h\|_\infty
        \left(
        e^{-L}F(L)+\int_0^L e^{-s}F(s)\d s
        \right)                        \\
         & \geq
        \|h\|_\infty
        \left(
        e^{-L}F(L)+\int_0^L e^{-s} \left(\frac{s}{L}\right)^d  F(L)\d s
        \right)                        \\
         & =
        \|h\|_\infty
        \frac{F(L)}{L^d}
        \left(
        L^de^{-L}+\int_0^L e^{-s} s^d\d s
        \right)                        \\
         & =
        \|h\|_\infty
        \frac{F(L)}{L^d}
        \left(
        d \int_0^L e^{-s} s^{d-1}\d s
        \right)\,,
    \end{align}
    and since $L>d$, the last integral is at least a numerical constant times
    $d!$. Thus
    \begin{align}
        \|h\|_\infty F(L)
        \le
        C\frac{L^d}{d!}.
        \label{eq:log_concave_sublevel_volume_factorial}
    \end{align}

    Next, for $s\ge L$, \eqref{eq:sublevel_volume_ratio_monotone} gives
    \begin{align}
        F(s)\le \left(\frac{s}{L}\right)^d F(L).
    \end{align}
    Therefore, again by integration by parts,
    \begin{align}
        \PP(V(Y)>L)
         & =
        \|h\|_\infty
        \int_L^\infty e^{-s}\d F(s)
        \le
        \|h\|_\infty
        \frac{F(L)}{L^d}
        d\int_L^\infty e^{-s}s^{d-1}\d s.
    \end{align}
    Combining this with \eqref{eq:log_concave_sublevel_volume_factorial},
    and using the exact formula for the Gamma tail,
    \begin{align}
        d\int_L^\infty e^{-s}s^{d-1}\d s
        =
        d!e^{-L}\sum_{j=0}^{d-1}\frac{L^j}{j!},
    \end{align}
    gives
    \begin{align}
        e^L\PP(V(Y)>L)
        \le
        C\sum_{j=0}^{d-1}\frac{L^j}{j!}.
        \label{eq:log_concave_potential_tail_factorial}
    \end{align}
    Finally,
    \begin{align}
        \E\left[
            \min\left\{e^{V(Y)},e^L\right\}
            \right]
         & =
        \E\left[
            e^{V(Y)}\mathbf{1}_{\{V(Y)\le L\}}
            \right]
        +
        e^L\PP(V(Y)>L)
        \\
         & =
        \|h\|_\infty F(L)
        +
        e^L\PP(V(Y)>L)
        \\
         & \le
        C
        \left(
        \frac{L^d}{d!}
        +
        \sum_{j=0}^{d-1}\frac{L^j}{j!}
        \right)
        =
        C
        \sum_{j=0}^d\frac{L^j}{j!}.
    \end{align}
    This proves the result.
\end{proof}


\begin{theorem}[Log-concave kernels pay \(d\) logarithms with factorial dependence]
    \label{theorem:log_concave_d_log_factor_factorial}
    Assume that $q$ is a log-concave probability density on $\Rd$.
    Then there exists a numerical constant $C>0$ such that for every
    $\lambda\in(0,1)$ and every $a\in(0,1]$,
    \begin{align}
        D_{f_a}(\PP_{X,T_0}\|\PP_X\otimes\PP_{T_0})
        \le
        \frac{C}{(1-\lambda)^d}
        a^2
        \sum_{j=0}^d
        \frac{1}{j!}
        \log^j\left(\frac{e}{a}\right).
        \label{eq:log_concave_d_log_factor_factorial_fa_upper}
    \end{align}
    In particular, for every $0<\eta\le \lambda$,
    \begin{align}
        \sup_{f:\K\to[0,1]}
        \Lambda_{\mathrm{var},q}(E_\lambda,p_X;f)
        \le
        \frac{C}{(1-\lambda)^d}
        \frac{\eta^2}{\lambda^2}
        \sum_{j=0}^d
        \frac{1}{j!}
        \log^j\left(\frac{e\lambda}{\eta}\right)
        +
        \frac{\eta^2}{2}.
        \label{eq:log_concave_d_log_factor_factorial_variance_upper}
    \end{align}
\end{theorem}

\begin{proof}
    Suppose that $d>1$, then
    \begin{align}
        B_\lambda\triangleq\frac{1}{(1-\lambda)^d}\,,
        \qquad
        \text{and when }\lambda=\frac{1}{d}, \qquad e < B_\lambda < 2e\,.
    \end{align}
    % \paragraph{Case 1: $aB_\lambda\ge1$.}
    % We have
    % \begin{align}
    %     \E\left[
    %         R_0\mathbf{1}_{\{R_0\le 1/a\}}
    %         \right]
    %     \le
    %     \frac1a,
    %     \qquad
    %     \PP\left(R_0>\frac1a\right)\le1.
    % \end{align}
    % Hence
    % \begin{align}
    %     D_{f_a}(\PP_{X,T_0}\|\PP_X\otimes\PP_{T_0})
    %     \le
    %     C a.
    % \end{align}
    % Since $aB_\lambda\ge1$, we have $a\le a^2B_\lambda$, so
    % \begin{align}
    %     D_{f_a}(\PP_{X,T_0}\|\PP_X\otimes\PP_{T_0})
    %     \le
    %     C a^2B_\lambda
    %     \le
    %     C a^2B_\lambda
    %     \sum_{j=0}^d
    %     \frac{1}{j!}
    %     \log^j\left(\frac{e}{a}\right).
    %     \label{eq:log_concave_d_factorial_large_a}
    % \end{align}

    % \paragraph{Case 2: $aB_\lambda<1$.}
    % Suppose that $aB_\lambda<1$, the other case is trivial since $a\le a^2B_\lambda$.
    Recall the self-identical decomposition
    \begin{align}
        X'=\lambda\sum_{k=1}^\infty (1-\lambda)^{k-1}T_k,
        \qquad
        X=\lambda T_0+(1-\lambda)X'\,.
    \end{align}
    Since $q$ is log-concave, $p$ is log-concave.
    Write
    \begin{align}
        p(x)=\|p\|_\infty e^{-V(x)}
    \end{align}
    with $V$ convex and $\inf_x V(x)=0$. From \eqref{eq:R0_identity_1},
    \begin{align}
        R_0
        =
        B_\lambda
        \frac{p(X')}{p(X)}
        \le
        B_\lambda
        \frac{\|p\|_\infty}{p(X)}
        =
        B_\lambda e^{V(X)}.
        \label{eq:log_concave_d_factorial_R0_upper}
    \end{align}

    We use the same quadratic-then-linear reduction as in the one-dimensional
    proof:
    \begin{align}
        D_{f_a}(\PP_{X,T_0}\|\PP_X\otimes\PP_{T_0})
        \le
        e^a
        a^2
        \left(
        1+
        \E\left[
            R_0\mathbf{1}_{\{R_0\le 1/a\}}
            \right]
        \right)
        +
        a\PP\left(R_0>\frac1a\right).
        \label{eq:log_concave_d_factorial_reduction}
    \end{align}
    Suppose that $aB_\lambda<1$ since the other case is trivial since $a\le a^2B_\lambda$, and define
    \begin{align}
        L=\log\left(\frac{1}{aB_\lambda}\right)>0.
    \end{align}
    From \eqref{eq:log_concave_d_factorial_R0_upper},
    \begin{align}
        \E\left[
            R_0\mathbf{1}_{\{R_0\le 1/a\}}
            \right]
         & \le
        \E\left[
            \min\left\{R_0,\frac1a\right\}
            \right]
        \le
        B_\lambda
        \E\left[
            \min\left\{e^{V(X)},e^L\right\}
            \right].
    \end{align}
    Applying \cref{lem:log_concave_truncated_exp_potential} with $h=p$ and
    $Y=X$ gives
    \begin{align}
        \E\left[
            R_0\mathbf{1}_{\{R_0\le 1/a\}}
            \right]
        \le
        C B_\lambda
        \sum_{j=0}^d \frac{L^j}{j!}.
        \label{eq:log_concave_d_factorial_truncated_mean}
    \end{align}
    Similarly,
    \begin{align}
        \PP\left(R_0>\frac1a\right)
         & \le
        \PP\left(B_\lambda e^{V(X)}>\frac1a\right)
        =
        \PP(V(X)>L).
    \end{align}
    Since
    \begin{align}
        e^L\PP(V(X)>L)
        \le
        \E\left[
            \min\left\{e^{V(X)},e^L\right\}
            \right],\label{ineq:F_L_is_high_probability_for_X}
    \end{align}
    another application of \cref{lem:log_concave_truncated_exp_potential}
    yields
    \begin{align}
        a\PP\left(R_0>\frac1a\right)
         & \le
        ae^{-L}
        C
        \sum_{j=0}^d \frac{L^j}{j!}
        =
        C a^2B_\lambda
        \sum_{j=0}^d \frac{L^j}{j!}.
        \label{eq:log_concave_d_factorial_tail_term}
    \end{align}
    Plugging \eqref{eq:log_concave_d_factorial_truncated_mean} and
    \eqref{eq:log_concave_d_factorial_tail_term} into
    \eqref{eq:log_concave_d_factorial_reduction}, and using
    $B_\lambda\ge1$, $a\le1$, and $L\le \log(e/a)$, gives
    \begin{align}
        D_{f_a}(\PP_{X,T_0}\|\PP_X\otimes\PP_{T_0})
         & \le
        C a^2B_\lambda
        \sum_{j=0}^d
        \frac{1}{j!}
        \log^j\left(\frac{e}{a}\right)\,.
    \end{align}
\end{proof}

\subsection{Showing log-concavity of $p_X$}
It remains to show that $p_X$ is log-concave, which is what we basically used in the proofs of the results in the previous section (actually we assumed log-concavity of $q$, which implies log-concavity of $p_X$, but in general we only need log-concavity of $p_X$).

We start by a general control over the function $r$ and $s_x$, which we defined in \cref{eq:s_in_terms_of_r} as
\marginpar{The $\bar{s}_x$ here is actually the $s_x - \mathbf{1}$ we defined before. This just has the convenient property that $\cE_x^\ell(t) = \ell(x)\bar{s}_x(t)$.}
\begin{align}
    r(x, t) & \triangleq  \frac{p_{X, T_0}(x, t)}{p_{T_0}(t) p_X(x)}\,, \quad\text{then}\quad
    \bar{s}_x(t) = \frac{r(x, t) - 1}{\lambda}\,.
\end{align}

Suppose that the core $p$ is log-concave, defined by
\begin{align}
    p(z) = \|p\|_\infty e^{-V(z)}\qquad \text{with $V:\R^d\to\R$ convex and $\inf_{z\in\K} V(z)=0$.}
\end{align}
Then, define
\begin{align}
    F_L = \left\{z \in \K: V(z) \le L\right\},
    \qquad
    \text{which is convex.}
\end{align}
We prove that for all $t \in F_L$,
\begin{align}
    r(x, t) \leq (1 - \lambda)^{-d} e^{\lambda L}.
\end{align}
To see this, recall the alternative definition of $r$ as
\begin{align}
    r(x, t) & = \frac{1}{(1-\lambda)^d} \cdot \frac{p\left(\frac{x-\lambda t}{1-\lambda}\right)}{p(x)}\,.
\end{align}
Then write
\begin{align}
    x    & = (1-\lambda) \frac{x - \lambda t}{1-\lambda} + \lambda t                       \\
    \implies\qquad
    V(x) & \leq (1-\lambda) V\left(\frac{x - \lambda t}{1-\lambda}\right) + \lambda V(t)   \\
    \implies\qquad
    p(x) & \geq p\left(\frac{x - \lambda t}{1-\lambda}\right)^{(1-\lambda)} p(t)^{\lambda} \\
    \implies\qquad
    \frac{p\left(\frac{x - \lambda t}{1-\lambda}\right)}{p(x)}
         & \leq  \left(\frac{p(x)}{p(t)}\right)^{\lambda/(1-\lambda)}
    \\
    \implies\qquad
    r(x, t)
         & \leq
    (1-\lambda)^{-d}
    \left(\frac{p(x)}{p(t)}\right)^{\lambda/(1-\lambda)}
    \approx
    C \exp\left( \frac{V(x) - V(t)}{d}\right)\,,
\end{align}
where the last approximation is for the choice of $\lambda = 1/d$.
By doing things slightly differently we can also get the sharper bound directly. Let
\[
    y=\frac{x-\lambda t}{1-\lambda}.
\]
Since \(x=(1-\lambda)y+\lambda t\), convexity gives
\[
    V(x)\le (1-\lambda)V(y)+\lambda V(t),
\]
and therefore
\[
    \frac{p(y)}{p(x)}
    =
    \exp\{V(x)-V(y)\}
    \le
    \exp\{\lambda(V(t)-V(y))\}
    \le e^{\lambda L},
\]
where we used \(V(y)\ge 0\) and \(t\in F_L\). Plugging this into the alternative formula for \(r\) gives
\[
    r(x,t)\le (1-\lambda)^{-d}e^{\lambda L}.
\]


So on the high-probability set $F_L$, we have the uniform bound
\begin{align}
    \sup_{t \in F_L} r(x, t) \leq (1 - \lambda)^{-d} e^{\lambda L}.
\end{align}
This gives us a uniform control over $s_x(t) = (r(x, t) - 1)/\lambda$ on the high-probability set $F_L$:
\begin{align}
    \sup_{t \in F_L} s_x(t) \leq \frac{(1 - \lambda)^{-d} e^{\lambda L} - 1}{\lambda}.
\end{align}
Also the set $F_L$ has high probability under $p$, as we have proved in \cref{ineq:F_L_is_high_probability_for_X} that
\begin{align}
    \PP(V(X)>L)
    \le
    e^{-L}
    \E\left[
        \min\left\{e^{V(X)},e^L\right\}
        \right]
    \le
    e^{-L}
    C
    \sum_{j=0}^d \frac{L^j}{j!}\,,
\end{align}